\section{CHALLENGES AND SOLUTIONS}
To address the limitations of conventional browser fuzzing for layout
invalidation correctness, we introduce \toolname: a system for
finding, minimizing, and grouping under-invalidation bugs in modern web
browser layout engines. In the following, we outline core challenges
that shape \toolname's design, along with their solutions aimed
at effectively fuzzing browser layout invalidation.

Existing browser fuzzers do not provide this combination of a targeted
state transition and a same-browser correctness oracle. JavaScript,
WebGL, WebAssembly, and origin fuzzers focus on different execution
surfaces and usually use crashes, security violations, or API failures as
their signal~\cite{han2019codealchemist,park2020die,peng2023gleefuzz,draissi2025wemby,kim2022fuzzorigin}.
Rendering fuzzers instead compare results across browser versions,
browsers, or executions~\cite{song2022r2z2,song2023metamong,zhou2025janus}.
Those techniques can reveal valuable regressions, but 
not identify whether an incremental layout result is stale.

\subsection{Challenge 1: Generating Inputs that Trigger Under-Invalidation Bugs}

The first challenge is generating inputs that are both valid web pages
and useful tests for layout invalidation. A browser fuzzer can easily
create HTML and CSS that the browser accepts, but such inputs may never
exercise the incremental layout paths where under-invalidation bugs
occur. To expose these bugs, the generated page must include an initial
layout, a meaningful DOM or CSS update, and a final state where the
browser should recompute affected geometry. \toolname therefore
biases generation toward CSS properties and mutations that are likely to
change layout, so that each test is more likely to trigger the browser's
invalidation logic rather than only testing parsing or general rendering
behavior. This is deliberately narrower than render-tree mutation and
cross-rendering approaches, which do not use the incremental-versus-fresh
layout transition as their test oracle~\cite{zeng2026rtfuzz,song2022r2z2,song2023metamong,zhou2025janus}.

\textbf{Display}: The \texttt{display} property is a useful target
because it determines the layout mode used for an element and
changes how its children participate in layout. Changing \texttt{display}
can move an element between block, inline, flex, grid, table, or hidden
states, each of which requires different layout behavior. These
changes can affect both the modified element and its descendants, they
provide strong opportunities to expose missed invalidation decisions in
incremental layout engines~\cite{panchekha2024reusing,kilpatrick2021layoutng}.

\textbf{Grid}: Grid layout creates many invalidation opportunities
because it introduces relationships between a container and
its children. Changing grid templates, gaps, placement rules, or
auto-flow can reposition multiple children from a single style update.
This makes grid a strong target for \toolname because a missed
invalidation can leave incorrect child positions or incorrect container
geometry. Such dependent geometry is precisely the kind of cached result
that incremental-layout implementations must invalidate correctly~\cite{panchekha2024reusing,liu2023medea}.

\textbf{Flow-root}: The \texttt{display:flow-root} value is important
because it establishes a new block formatting context and changes the
boundary at which surrounding layout is resolved. Mutating it can therefore
change both the target box and the geometry of neighboring content, making
it useful for exercising invalidation propagation~\cite{panchekha2024reusing}.

\textbf{Width and height}: Width and height changes directly affect an
element's box geometry. When an element changes size, the browser may
need to recompute that element, its descendants, siblings, and sometimes
ancestors. These mutations are useful because layout engines often
appears as an obvious mismatch in element size or position.

\textbf{Percentage and pixel sizes}: Percentage sizes test
parent-dependent layout behavior because the final size of an element
depends on its containing block. Pixel sizes create more direct change
to the geometry possibly causing conflict issues with surrounding 
elements. Using both lets \toolname exercise
dependent layout constraints to produce failures that are
easy to observe and reproduce. General browser fuzzers need not preserve
such a parent-dependent update relationship, because their target oracles
are typically crashes, security violations, or externally observed output
differences~\cite{han2019codealchemist,park2020die,kim2022fuzzorigin,song2022r2z2}.

\textbf{Transforms}: Transform changes affect visual position and
reported geometry. They provide a controlled update whose effect can be
checked through the geometry oracle while interacting with containing blocks
and descendant coordinates. This makes them useful for distinguishing a
stale incremental result from a freshly computed final layout.

\textbf{Noise reduction}: \toolname disables properties such as
\texttt{writing-mode}, \texttt{content-visibility}, padding, and table
display values because they are less likely to produce
under-invalidation bugs. Instead, these properties
often add unnecessary complexity to the generated tests. Removing them 
keeps the generated tests focused on mutations
that are more likely to expose missed invalidation decisions and helps
produce smaller, clearer bug reports. This focus is important because
unconstrained random inputs can make an observed rendering difference hard
to attribute to one invalidation decision~\cite{song2022r2z2,song2023metamong,zhou2025janus}.

\textbf{Solution 1}: \toolname includes a generation strategy for
steering generated inputs toward page updates that are likely to trigger
layout under-invalidation bugs. The system focuses on CSS properties and
mutations that should force layout recomputation, increasing the chance
that a missed invalidation decision becomes observable.

\subsection{Challenge 2: Minimizing and Isolating Under-Invalidation Bugs}

Under-invalidation bugs appear in
large generated pages that contain many elements, styles, and mutations.
Even when only a small part of the page is responsible for the failure,
the original test case may include many unrelated structures introduced
by random generation. Such inputs are not useful for diagnosis because
they obscure the DOM structure or CSS change that triggered the stale
layout state.

This differs from reducing a conventional crash or a cross-browser output
discrepancy: the reducer must preserve not only a failing final page but
also the particular incremental state transition that exposes the stale
result~\cite{zellerFuzzingBookReducer,misherghi2006hdd,webkitTestCaseReduction}.

\textbf{Preserving the bug trigger}: Minimization is difficult because
under-invalidation bugs depend on a precise sequence of layout states.
The browser first lays out the original page, then applies a DOM or CSS
update, and then decides which cached layout information can be reused.
If the minimizer removes an element, style, or mutation that seems
unrelated but changes this invalidation path, the bug may disappear even
though the original conditions were present.

\textbf{Separating base and modified styles}: To identify which declarations
are essential to the transition, \toolname moves one modified style
declaration at a time into the base style map. If the mismatch persists after
the move, that declaration was not required as an incremental update and can
be removed from the modified state. If the mismatch disappears, the
declaration remains modified because changing it after the initial layout is
part of the bug trigger.

\textbf{Solution 2}: \toolname reduces bug-triggering inputs to their 
essential elements, isolates the structures that cause the failure, and produces
a test case that is easier to analyze and debug. Its reduction predicate is
the same-browser incremental-versus-fresh geometry mismatch.

\subsection{Challenge 3: Grouping and Categorizing Under-Invalidation Bugs}

A major challenge in layout fuzzing is that
many generated inputs can trigger the same underlying browser bug. Two
test cases may differ slightly in the HTML or CSS while still
exercising the same missed invalidation decision inside the layout
engine. Without grouping, developers would need to inspect many
redundant reports before finding the distinct bugs that actually require
attention.

Prior rendering fuzzers can produce many distinct-looking test cases from
their mutation and comparison strategies, but their output comparisons do
not identify the shared invalidation structure inside one layout
engine~\cite{song2022r2z2,song2023metamong,zhou2025janus}. Consequently,
deduplication for this bug class must retain the DOM and style transition,
rather than grouping only by a rendered image or a failure signal.

\textbf{Structural representation}: Grouping under-invalidation bugs
requires representing the parts of a test case that matter for the
failure. The system must capture the DOM structure, base styles,
modified styles, and the element relationships involved for each
layout state. A consistent representation makes it possible to compare
bugs by their underlying layout pattern rather than by superficial
differences in generated HTML.

\textbf{Automated categorization}: Categorization is difficult because
the process must be automated while still avoiding overly broad groups.
If the grouping rule is too strict, the system leaves many duplicates
unmerged; if it is too loose, unrelated bugs may be combined into the
same category. \toolname therefore needs grouping logic that can
recognize repeated bug patterns while preserving enough detail to keep
distinct failures separate.

\textbf{Skipping minimization}: Full minimization repeatedly
executes a candidate in the browser and is expensive when a newly detected
failure already matches a known under-invalidation pattern. \toolname first
compares a candidate with existing structural group rules. If the candidate
matches a group, the system saves it as another instance of that group and
skips full minimization. The minimized representative and shared rule already
provide the information needed for triage; unmatched candidates are still
minimized before they are saved or used to form a new group.

\textbf{Solution 3}: \toolname groups bugs using structural rules
and merged bug representations. This reduces duplicate reports, keeps
similar failures together, and lets developers focus on the distinct
under-invalidation patterns that are most likely to be actionable.
